\begin{theorem}[Clifton--Hogarth]\label{thm:flat}
For every becoming theory, $B_\Mink\in\{\Ichr,\Jm,\U\}$.
\end{theorem}

\begin{proof}
Proper orthochronous Poincaré maps are automorphisms of $\Mink$ in $\GH$, so
$B_\Mink$ is invariant under them by (LC). Clifton and Hogarth's Theorem~1
\citep[p.~365]{CH1995} gives the conclusion from this invariance,
transitivity and chronological becoming. Its proof (pp.~365--367) uses only
point-preserving spatial rotations, boosts and translations, and they note
that it does not use spatial reflections (pp.~385--386). Their stated
hypothesis is invariance under all time-orientation-preserving isometries,
reflections included; using only the proper orthochronous subgroup is a
weakening whose adequacy rests on their proof, not their statement.
\end{proof}

\begin{theorem}[No becoming outside the causal past]\label{thm:main}
Let $B$ be a becoming theory.
\begin{enumerate}
\item[(a)] If $B_\Mink\neq\U$, then $B_M\subseteq\Jm_M$ for every $M\in\GH$.
\item[(b)] If $B_\Mink=\U$, then $B_M=\U_M$ for every $M\in\GH$.
\end{enumerate}
Hence either $B$ is universal everywhere, or $\Ichr_M\subseteq
B_M\subseteq\Jm_M$ for every $M$.
\end{theorem}

Call $B$ \emph{universal} if $B_M=\U_M$ for every $M$. By
Theorem~\ref{thm:main} this holds if and only if $B_\Mink=\U$, and if $B_M=\U_M$
for one $M$ then it holds for all; otherwise $B$ is \emph{non-universal}.

\begin{corollary}[No shared present]\label{cor:nopresent}
Call $x$ and $b$ \emph{co-present} in $M$ if $x\,B_M\,b$ and $b\,B_M\,x$. For
every non-universal becoming theory, co-presence is the identity relation in
every $M\in\GH$. Equivalently, no becoming theory is local, non-trivial and
admits two distinct co-present events in some globally hyperbolic spacetime.
\end{corollary}

\begin{proof}
By Theorem~\ref{thm:main}(a), $x\le b$ and $b\le x$, and globally hyperbolic
spacetimes are causal, so $x=b$.
\end{proof}

\begin{corollary}[Worldline-dependent becoming]\label{cor:worldline}
Suppose that for each $M\in\GH$, each inextendible timelike curve $\lambda$ in
$M$ and each $b\in\lambda$ a set $B_{M,\lambda}(b)$ is given, satisfying
Clifton--Hogarth's worldline becoming ($\Ichr(b)\cap\lambda\subseteq
B_{M,\lambda}(b)$) and worldline transitivity ($a\in B_{M,\lambda}(b)$ and
$b\in B_{M,\lambda'}(c)$ imply $a\in B_{M,\lambda'}(c)$) \citep[pp.~362--363]{CH1995},
and local covariance in the form: for a morphism $\psi:N\to M$, an inextendible
timelike $\lambda_N$ in $N$ and any inextendible timelike $\lambda$ in $M$
extending $\psi\circ\lambda_N$, $a\in B_{N,\lambda_N}(b)\iff\psi(a)\in
B_{M,\lambda}(\psi(b))$. Then $B_{M,\lambda}(b)$ does not depend on $\lambda$, and
the resulting relation is a becoming theory. Hence Theorem~\ref{thm:main}
applies to worldline-dependent becoming.
\end{corollary}

Clifton and Hogarth anticipated that extending their Theorem~2 would be
harder still (``much less to extend Theorem 2'', p.~385). Under local
covariance it is not: as they themselves note (p.~363), worldline transitivity
makes observers whose worldlines cross agree, and the remaining work is done by
Theorem~\ref{thm:main}.

Write $E^-=\Jm\setminus\Ichr$ for the horismos (lightlike-only relation).

\begin{theorem}[Classification of the residual freedom]\label{thm:class}
Let $B$ be a non-universal becoming theory and $N_M=B_M\cap E^-_M$. Then $N$
is locally covariant and satisfies
\begin{enumerate}
\item[(NC)] if $(x,y),(y,b)\in N_M$ and $(x,b)\in E^-_M$, then $(x,b)\in N_M$.
\end{enumerate}
Conversely, if $N\subseteq E^-$ is locally covariant and satisfies (NC), then
$B=\Ichr\cup N$ is a becoming theory. On $\Mink$, $N$ is empty or all of $E^-$.
\end{theorem}

\begin{theorem}[Non-uniqueness]\label{thm:weyl}
Let $N_M(x,b)$ hold iff $x\in E^-(b)$, $x\neq b$, and some lightlike geodesic
segment from $x$ to $b$ contains a point at which the Weyl tensor is nonzero.
Then $\Ichr\cup N$ is a becoming theory equal to $\Ichr$ on $\Mink$, equal to
$\Jm$ on any spacetime with nowhere-vanishing Weyl tensor (such as the
Schwarzschild exterior), and strictly between $\Ichr$ and $\Jm$ on spacetimes
that are flat on a nonempty causally convex open set and have nonzero Weyl
tensor somewhere. It depends only on the conformal structure (not merely the
curvature: conformally flat spacetimes such as FLRW give $\Ichr$).
\end{theorem}

\begin{proposition}[Regularity]\label{prop:reg}
Let $B$ be non-universal. If every $B_M$ is closed in $M\times M$, then
$B=\Jm$. If every $B_M\setminus\Delta_M$ is open, then $B=\Ichr$.
\end{proposition}
